% !TeX root = ../Book1.tex
\begin{figure}[htbp]
	\centering
	\begin{tikzpicture}[scale=1,>=stealth,line cap=round,line join=round]
		\begin{scope}[xshift=-3cm]
			\draw[thick] (0,0) ellipse (1.6 and 1.05);
			\draw[->,blue!70!black,thick] (1.6,0) -- (1.6,0.72) node[right] {\(p\)};
			\draw[->,blue!70!black,thick] (0,1.05) -- (-0.72,1.05) node[above] {\(v\)};
			\draw[->,blue!70!black,thick] (-1.6,0) -- (-1.6,-0.72) node[left] {\(-p\)};
			\draw[->,blue!70!black,thick] (0,-1.05) -- (0.72,-1.05) node[below] {\(-v\)};
			\node at (0,-1.85) {平面凸曲线};
		\end{scope}
		\draw[->,thick] (-0.75,0) -- (0.75,0) node[midway,above] {\(T\)};
		\begin{scope}[xshift=3cm]
			\draw[thick,blue!70!black] (0,0) circle (1.15);
			\draw[->,thick,blue!70!black] (80:1.15) arc (80:120:1.15);
			\draw[->,thick,blue!70!black] (260:1.15) arc (260:300:1.15);
			\fill (0,1.15) circle (1.5pt) node[above] {\(p\), \(u=0,2\pi\)};
			\fill (-1.15,0) circle (1.5pt) node[left] {\(v\)};
			\fill (0,-1.15) circle (1.5pt) node[below] {\(-p\), \(u=\pi\)};
			\fill (1.15,0) circle (1.5pt) node[right] {\(-v\)};
			\node at (0,-1.85) {球面中的一个大圆};
		\end{scope}
	\end{tikzpicture}
	\caption{等号情形中，切向示性曲线按累计曲率 \(u\) 沿大圆前进一周。这里以椭圆示意原曲线；一般凸曲线若有直线段，切向量会在相应方向停顿。}
	\label{g1:fig:total-curvature}
\end{figure}
